When a language model judges a piece of text as gripping, overwrought, or merely flat, it makes an affect-laden subjective call, and different readers will not agree on it. The distinction between drama and \emph{melodrama} sharpens the question: is melodrama simply more drama, or something categorically different, and does reader disagreement live on the same axis? No prior work studies how models represent drama; the broader debate over whether subjective constructs are single directions or multi-dimensional subspaces remains unsettled, and model-internal affect geometry has never been joined to per-reader behavioral variation. We probe the representation geometry of Qwen2.5-7B and 1.5B using a controlled, blind-validated drama corpus and 70 identified GoEmotions raters, with difference-in-means directions, participation-ratio estimates of subspace dimension, and nested held-out logistic models gauged against a shuffled-reader null. We find a clean dissociation. The construct is multi-dimensional: drama decodes at AUROC 1.00, yet the step to melodrama rotates to roughly 83 degrees off the drama axis (subspace participation ratio about 5.6), so melodrama lives on a distinct excess direction. Reader variation, by contrast, is a sliding threshold: adding a per-reader threshold lifts held-out AUROC from 0.51 to 0.62 and 0.66 to 0.70, while letting each reader pick their own direction adds exactly the shuffled-reader null, and six prompted personas leave the drama axis fixed at cosine at least 0.97. Different kinds of drama live in shared content; different people mostly move a threshold.
